\documentclass[10pt,a4paper]{amsart}
\usepackage{amsmath,amssymb,amsthm,graphicx}
% Retained preamble for the Stacks Project chapter snapshot in this directory.
% The chapter's own preamble.tex is a copy of the shared Stacks Project preamble; the block
% below is sliced MECHANICALLY, line for line, from it. Nothing was added, removed or
% redefined; the wrapper class, the xr-hyper/hyperref wiring and the externaldocument list
% are omitted because the assembler supplies the document class.
%
% lines taken from preamble.tex: 3, 5, 9, 11, 12, 156, 157, 158, ... (51 lines total)
%
% package coverage restored from the chapter's own preamble (2026-10-13): xy, xy-2cell

\usepackage[all]{xy}
\xyoption{2cell}
\UseAllTwocells

\newenvironment{reference}{\comment}{\endcomment}
\newenvironment{slogan}{\comment}{\endcomment}
\newenvironment{history}{\comment}{\endcomment}
\theoremstyle{plain}
\newtheorem{theorem}[subsection]{Theorem}
\newtheorem{proposition}[subsection]{Proposition}
\newtheorem{lemma}[subsection]{Lemma}
\theoremstyle{definition}
\newtheorem{definition}[subsection]{Definition}
\newtheorem{example}[subsection]{Example}
\newtheorem{exercise}[subsection]{Exercise}
\newtheorem{situation}[subsection]{Situation}
\theoremstyle{remark}
\newtheorem{remark}[subsection]{Remark}
\newtheorem{remarks}[subsection]{Remarks}
\numberwithin{equation}{subsection}
\def\lim{\mathop{\mathrm{lim}}\nolimits}
\def\colim{\mathop{\mathrm{colim}}\nolimits}
\def\Spec{\mathop{\mathrm{Spec}}}
\def\Hom{\mathop{\mathrm{Hom}}\nolimits}
\def\Ext{\mathop{\mathrm{Ext}}\nolimits}
\def\SheafHom{\mathop{\mathcal{H}\!\mathit{om}}\nolimits}
\def\SheafExt{\mathop{\mathcal{E}\!\mathit{xt}}\nolimits}
\def\Sch{\mathit{Sch}}
\def\Mor{\mathop{\mathrm{Mor}}\nolimits}
\def\Ob{\mathop{\mathrm{Ob}}\nolimits}
\def\Sh{\mathop{\mathit{Sh}}\nolimits}
\def\NL{\mathop{N\!L}\nolimits}
\def\CH{\mathop{\mathrm{CH}}\nolimits}
\def\proetale{{pro\text{-}\acute{e}tale}}
\def\etale{{\acute{e}tale}}
\def\QCoh{\mathit{QCoh}}
\def\Ker{\mathop{\mathrm{Ker}}}
\def\Im{\mathop{\mathrm{Im}}}
\def\Coker{\mathop{\mathrm{Coker}}}
\def\Coim{\mathop{\mathrm{Coim}}}
\DeclareMathSymbol{\boxtimes}{\mathbin}{AMSa}{"02}
\def\QCohstack{\mathcal{QC}\!\mathit{oh}}
\def\Cohstack{\mathcal{C}\!\mathit{oh}}
\def\Spacesstack{\mathcal{S}\!\mathit{paces}}
\def\Quotfunctor{\mathrm{Quot}}
\def\Hilbfunctor{\mathrm{Hilb}}
\def\Curvesstack{\mathcal{C}\!\mathit{urves}}
\def\Polarizedstack{\mathcal{P}\!\mathit{olarized}}
\def\Complexesstack{\mathcal{C}\!\mathit{omplexes}}
\def\Pic{\mathop{\mathrm{Pic}}\nolimits}
\def\Picardstack{\mathcal{P}\!\mathit{ic}}
\def\Picardfunctor{\mathrm{Pic}}
\def\Deformationcategory{\mathcal{D}\!\mathit{ef}}
\usepackage{cleveref}
\crefname{originaldefinition}{original definition}{original definitions}
\Crefname{originaldefinition}{Original definition}{Original definitions}
\makeatletter
\newlabel{definition-big-fppf-site}{{021R}{source}{}{sourcepointer-1}{}}
\newlabel{definition-big-fppf-site@cref}{{[originaldefinition][0][]021R}{[1][0][]source}{}{}{}}
\makeatother
\crefname{originallemma}{original lemma}{original lemmas}
\Crefname{originallemma}{Original lemma}{Original lemmas}
\makeatletter
\newlabel{lemma-fppf}{{021O}{source}{}{sourcepointer-2}{}}
\newlabel{lemma-fppf@cref}{{[originallemma][0][]021O}{[1][0][]source}{}{}{}}
\makeatother
\crefname{originalreference}{original reference}{original references}
\Crefname{originalreference}{Original reference}{Original references}
\makeatletter
\newlabel{sets-lemma-what-is-in-it}{{000R}{source}{}{sourcepointer-3}{}}
\newlabel{sets-lemma-what-is-in-it@cref}{{[originalreference][0][]000R}{[1][0][]source}{}{}{}}
\makeatother
\crefname{originalreference}{original reference}{original references}
\Crefname{originalreference}{Original reference}{Original references}
\makeatletter
\newlabel{sets-lemma-coverings-site}{{000X}{source}{}{sourcepointer-4}{}}
\newlabel{sets-lemma-coverings-site@cref}{{[originalreference][0][]000X}{[1][0][]source}{}{}{}}
\makeatother
\title{The induced topology on a subcategory from the big fppf topology}
\author{The Stacks Project Authors}
\date{stacks-project:topologies.tex:local-snapshot-20261012; GNU FDL 1.2 or later}
\begin{document}
\maketitle
\noindent\textit{Editorial note.} All delivered mathematical text is the author's own, copied verbatim from the pinned Stacks Project chapter file topologies.tex; no statement or proof step was rewritten, reordered or removed.\par\medskip
\noindent\textit{Editorial note.} Span recomputation by label pairing (LANE\_BRIEF 10.2): the official pool candidate recorded statement [2241, 2255] and proof(s) [[2257, 2283]]; this spec covers the complete environment instead (statement 2241-2255, proof span(s) [(2257, 2283)]), since in this chapter the author's proof environment(s) follow the statement environment and some pool entries under-counted them.\par\medskip
\ifdefined\hypertarget\hypertarget{sourcepointer-1}{}\fi
\noindent\textit{Source reference.} Original definition 021R: sources/stacks-chow/tags.txt, source lines 2680--2680. Cross-chapter standard prerequisite cited by the author inside the counted span; it is used as a pointer and its own proof is not part of this item.\par\medskip
\ifdefined\hypertarget\hypertarget{sourcepointer-2}{}\fi
\noindent\textit{Source reference.} Original lemma 021O: sources/stacks-chow/tags.txt, source lines 2677--2677. Cross-chapter standard prerequisite cited by the author inside the counted span; it is used as a pointer and its own proof is not part of this item.\par\medskip
\ifdefined\hypertarget\hypertarget{sourcepointer-3}{}\fi
\noindent\textit{Source reference.} Original reference 000R: sources/stacks-chow/tags.txt, source lines 36--36. Cross-chapter standard prerequisite cited by the author inside the counted span; it is used as a pointer and its own proof is not part of this item.\par\medskip
\ifdefined\hypertarget\hypertarget{sourcepointer-4}{}\fi
\noindent\textit{Source reference.} Original reference 000X: sources/stacks-chow/tags.txt, source lines 42--42. Cross-chapter standard prerequisite cited by the author inside the counted span; it is used as a pointer and its own proof is not part of this item.\par\medskip
\begin{lemma}
\label{lemma-fppf-induced}
Let $\Sch_{fppf}$ be a big fppf site as in
Definition \ref{definition-big-fppf-site}.
Let $T \in \Ob(\Sch_{fppf})$.
Let $\{T_i \to T\}_{i \in I}$ be an arbitrary fppf covering of $T$.
\begin{enumerate}
\item There exists a covering $\{U_j \to T\}_{j \in J}$ of $T$ in the site
$\Sch_{fppf}$ which refines $\{T_i \to T\}_{i \in I}$.
\item If $\{T_i \to T\}_{i \in I}$ is a standard fppf covering, then
it is tautologically equivalent to a covering of $\Sch_{fppf}$.
\item If $\{T_i \to T\}_{i \in I}$ is a Zariski covering, then
it is tautologically equivalent to a covering of $\Sch_{fppf}$.
\end{enumerate}
\end{lemma}

\begin{proof}
For each $i$ choose an affine open covering $T_i = \bigcup_{j \in J_i} T_{ij}$
such that each $T_{ij}$ maps into an affine open subscheme of $T$. By
Lemma \ref{lemma-fppf}
the refinement $\{T_{ij} \to T\}_{i \in I, j \in J_i}$ is an fppf covering
of $T$ as well. Hence we may assume each $T_i$ is affine, and maps into
an affine open $W_i$ of $T$. Applying
Sets, Lemma \ref{sets-lemma-what-is-in-it}
we see that $W_i$ is isomorphic to an object of $\Sch_{fppf}$.
But then $T_i$ as a finite type scheme over $W_i$
is isomorphic to an object $V_i$ of $\Sch_{fppf}$ by a second
application of
Sets, Lemma \ref{sets-lemma-what-is-in-it}.
The covering $\{V_i \to T\}_{i \in I}$ refines $\{T_i \to T\}_{i \in I}$
(because they are isomorphic).
Moreover, $\{V_i \to T\}_{i \in I}$ is combinatorially equivalent to a
covering $\{U_j \to T\}_{j \in J}$ of $T$ in the site
$\Sch_{fppf}$ by
Sets, Lemma \ref{sets-lemma-what-is-in-it}.
The covering $\{U_j \to T\}_{j \in J}$ is a refinement as in (1).
In the situation of (2), (3) each of the
schemes $T_i$ is isomorphic to an object of $\Sch_{fppf}$ by
Sets, Lemma \ref{sets-lemma-what-is-in-it},
and another application of
Sets, Lemma \ref{sets-lemma-coverings-site}
gives what we want.
\end{proof}

\end{document}
